\section{Pilot Unified Mobile Foundation}
\label{sec:pilot-unified-mobile-foundation}

The Pilot mobile product is the common private client for Personal Pilot, bounded
professional Workspaces, and the AION Boardroom.  It does not create three separate
accounts or three separate messaging systems.  One verified person uses one trusted
device and deliberately enters a governed space.  The selected space determines which
memory, conversations, tools, services, and approval powers are available.

This section records the implemented contract foundation and the first responsive
reference client.  The reference client is intentionally fixture-driven: it proves the
information architecture and boundary model without claiming that demonstration buttons
already possess live authority.  Live discovery, pairing, synchronization, messaging,
and action execution are connected in subsequent stages.

\subsection{Product Model}

The application exposes three primary modes:

\begin{description}
  \item[Personal.] Private conversation, inbox, tasks, calendar, household devices,
  television, games, learning, shopping, files, reminders, and personal approvals.
  \item[Workspace.] A deliberately bounded role in a client, employer, project, or
  professional engagement.  Its tools are supplied by an authorized workspace manifest.
  \item[Boardroom.] Mobile briefing, review, delegation, intervention, and approval for
  an owned or administered organization.  Complex business creation and administration
  remain desktop-first.
\end{description}

Changing the visible mode is not an authority grant.  Every request carries the active
person, device, space, membership, lease, and mother-brain context.  The receiving mother
brain validates that complete context again before reading private information or
executing an operation.

\subsection{Versioned Shared Contracts}

The canonical schema is version \texttt{pilot.unified.v1}.  Its implementation lives in
an isolated Python contract package.  It does not overwrite the existing AION cognitive
runtime, Pilot Fabric, GlyphNet, or Boardroom stores.

The first contract set defines:

\begin{itemize}
  \item person and persona references;
  \item trusted device and signed possession proof;
  \item Personal, Workspace, and Boardroom spaces;
  \item memberships, invitations, roles, and explicit scopes;
  \item short-lived, device-bound capability leases;
  \item the complete active context attached to every operation;
  \item participants, conversations, attachment references, and message envelopes;
  \item action proposals, approvals, executions, and immutable receipts;
  \item surface manifests, service-connection references, and memory references.
\end{itemize}

Every contract supports deterministic serialization and hashing.  Identifiers, timestamps,
and payload sizes are bounded before use.  Action states follow an explicit transition
graph so a client cannot convert a prepared proposal directly into a verified result.
Rendered controls confer no authority by themselves.

\subsection{Authority and Confused-Deputy Protection}

The validator treats the following tuple as one security boundary:

\begin{center}
\texttt{person + device + possession + space + membership + lease + mother}
\end{center}

The persona on the possession proof must match the active person.  The device on the
proof must match the requesting device.  The membership and capability lease must both
refer to the selected space and persona.  Their time windows and revisions must remain
valid, and the requested scope must be present in both records.  A bounded replay guard
rejects a nonce that has already been consumed.

This prevents a client from rendering a Boardroom button while retaining a Personal
lease, borrowing a scope from another Workspace, replaying an old approval, or persuading
one mother brain to exercise authority issued by another.

\subsection{Compatibility Without Destructive Migration}

Read-only adapters project existing stores into the new contracts.  They currently cover:

\begin{itemize}
  \item Pilot Fabric private identities and trusted phone devices;
  \item personal spaces, inbox tasks, messages, and contacts;
  \item governed communication drafts;
  \item GlyphNet thread events and durable conversation history;
  \item Boardroom workspace bundles and organizational context.
\end{itemize}

Existing records are not irreversibly rewritten.  A legacy approval without production
possession evidence is deliberately returned to an awaiting-approval state.  A Boardroom
bundle identifies a potential Workspace but grants no mobile capability until the
Boardroom issues a signed surface manifest and role lease.

\subsection{Ownership Registry}

The implementation registry prevents parallel subsystems from becoming competing sources
of truth.  Pilot Fabric owns private identity, device possession, personal tasks, service
connections, and television continuity.  GlyphNet owns conversation transport and its
thread journal.  The Boardroom owns organization records, departments, decisions, and
business policy.  GlyphChain proof modules may anchor hashes and receipts, but never
authorize an action merely because a chain entry exists.

Economic experiments---including PHO, PhotonPay, wallets, tokens, bonds, minting, and
staking---remain preserved but disabled in the mobile feature policy.  They are excluded
from routes, manifests, shortcuts, capability requests, and default navigation.  Radio
and wave transports remain optional adapters behind the same message envelope.

\subsection{Responsive Reference Client}

The first client is implemented inside the existing Next.js application at
\texttt{/pilot/mobile}.  It reads the same JSON fixture used by the contract tests, rather
than duplicating example identities and permissions in the interface.

The visible structure is intentionally close to a familiar messaging application:

\begin{enumerate}
  \item a compact Pilot header with private notification and person controls;
  \item a persistent identity strip showing person, space, role, mother, and connection;
  \item a Personal, Workspace, and Boardroom segmented switcher;
  \item horizontally scrollable, context-specific shortcuts from the surface manifest;
  \item a small stream of messages, tasks, reminders, decisions, approvals, and receipts;
  \item one persistent text and hold-to-speak Pilot composer;
  \item persistent Pilot, Inbox, Actions, Spaces, and Me navigation.
\end{enumerate}

{\setlength{\emergencystretch}{2em}
Personal shortcuts include Calendar, Tasks, Devices, TV, Games, Learning, Shopping, and
Files.  Workspace shortcuts are supplied by its bounded manifest.  Boardroom shortcuts
include Board, Sales, Marketing, Finance, Operations, Support, People, Products, and
Files.  They appear only when authorized.  The universal television and camera observer
entry points remain easy to reach.  Specialist functions no longer form one long
controller page.\par}

The three modes use distinct accent systems while retaining one component language.
Touch targets, safe-area insets, visible keyboard focus, screen-reader labels, scalable
type, contrast, and reduced-motion preferences are supported.  Protected notification
previews hide sender and content until the device proves possession and the user unlocks
the relevant space.

\subsection{Installable and Offline Shell}

The reference client includes a scoped web-application manifest, a Pilot aircraft icon,
and a network-first service worker.  The worker caches only the Pilot shell and same-origin
resources.  It attempts the current mother-brain response first and falls back to the
cached shell when connectivity is lost.  Offline rendering never implies that an action
was accepted, sent, or executed; such claims require a fresh signed receipt.

Production installation still requires trusted HTTPS, certificate and key lifecycle
management, safe remote discovery, and a native signed-phone path for biometric approval.
Those controls belong to the pairing and remote-continuity stage rather than this visual
foundation.

\subsection{Implemented State Language}

Cards use explicit states including \texttt{prepared}, \texttt{pending},
\texttt{locked}, \texttt{offline}, \texttt{expired}, \texttt{failed}, and
\texttt{verified}.  A prepared item is not approved; an approved item is not executed;
an executed item is not verified until an independent adapter or provider receipt proves
the outcome.  This language is shared across Personal, Workspace, and Boardroom surfaces.

\subsection{Verification Evidence and Honest Boundary}

Contract, compatibility, authority, replay, transition, economic-feature isolation,
manifest, installability, accessibility, and fixture-consumption checks are automated.
The complete AION Fabric suite contains 427 passing tests after this stage.  The local
reference route, manifest, and service worker each return a successful response.

The complete optimized Next.js production build compiles successfully and statically
generates the Pilot mobile route.  Live usability tests on representative iPhone and
Android sizes, replacement of the legacy controller stream, and real mother-brain
synchronization remain open and are not represented as complete.

\subsection{Next Implementation Boundary}

The next stage turns the phone into a secure client of a user-selected mother brain.  It
adds local visual pairing, signed possession, certificate pinning, local discovery,
remote relay or private-network continuity, explicit mother selection, revocation, and
recovery.  Only after that boundary is verified should the fixture cards become live
GlyphNet conversations, Boardroom reviews, Pilot actions, and external service approvals.
